% IMPORTANT: PLEASE USE XeLaTeX FOR TYPESETTING
\documentclass{sintefbeamer}

% packages, font, color, and newcommands
\usepackage{amsfonts, amsmath, oldgerm, lmodern}
\usepackage{xeCJK}
\usefonttheme{structuresmallcapsserif}

% meta-data
\title{Opmaak figuren}
\subtitle{Powerpoint en Inkscape}
\author{\href{j.busink@vu.nl}{}}
\date{\today}
\titlebackground{images/Template/VuBG2.png}
\setbeamertemplate{sidebar right}{}
\setbeamertemplate{footline}{%
\hfill\usebeamertemplate***{navigation symbols}
\hspace{1cm}\insertframenumber{}/\inserttotalframenumber}
\begin{document}

\maketitle





\section{Section 1}

% \begin{frame}{Beamer slides}{\thesection \, \secname}


\begin{frame}[fragile]{Software}{\thesection \, \secname}
      \framesubtitle{Subtitle}
      \begin{itemize}
            \item \hrefcol{mailto:j.busink@svu.nl}{Email: J.Busink} 
      \end{itemize}

      \begin{itemize}[<+->]
            \item \hrefcol{https://inkscape.org/release/inkscape-1.2.1/}{Inkscape (Windows, macOS, GNU/Linux)}
            \begin{figure}
                  \centering
                  \includegraphics[width = 2 cm]{images/Section 1/Inkscape_Logo.svg.png}
            \end{figure}
            \item Powerpoint
            \item Adobe Photoshop
      \end{itemize}
\end{frame}
% \end{frame}

\section{Section 2}

\begin{frame}[fragile]{Selecting the Class}
After the last update to the graphic profile, the \texttt{sintef} theme for
Beamer has been updated into a full-fledged class.
To start working with \texttt{sintefbeamer}, start a \LaTeX\ document with the
preamble:
\begin{block}{Minimum SINTEF Beamer Document}
\begin{lstlisting}[language=TeX]
\documentclass{sintefbeamer}
\begin{document}
\begin{frame}{Hello, world!}
\end {frame}
\end{lstlisting}
\end{block}
\end{frame}

\begin{frame}[fragile]{Simple Slide}
\framesubtitle{Subtitle}
\begin{itemize}[<+->]
\item Bullet 1
\item Bullet 2
\item Bullet 3
\end{itemize}
\begin{block}{Code for a Page with an Itemised List}<+->
\begin{lstlisting}[language=TeX]
\begin{frame}
  \frametitle{Writing a Simple Slide}
  \framesubtitle{It's really easy!}
  \begin{itemize}[<+->]
    \item A typical slide has bulleted lists
    \item These can be uncovered in sequence
  \end{itemize}
\end{frame}\end{lstlisting}
\end{block}
\end{frame}

\begin{frame}[fragile]{Look}
\begin{itemize}
\item To change the colour of the title dash, give one of the class options
      \texttt{cyandash} (default), \texttt{greendash}, \texttt{magentadash},
      \texttt{yellowdash}, or \texttt{nodash}.
\item To change between the light and dark themes, give the class options
      \texttt{light} (default) or \texttt{dark}. It is not possible to switch
      theme for one slide because of the design of Beamer---and it's probably a
      good thing.
\item To insert a final slide, use \verb|\backmatter|.
\item The aspect ratio defaults to 16:9, but you can change it to 4:3 for old
      projectors by passing the class option \texttt{aspectratio=43}; any other
      values accepted by Beamer are also possible.
\end{itemize}
\end{frame}

\section{Section 3}

\begin{frame}
\frametitle{Frame Title}
\begin{itemize}
\item Some text
\end{itemize}
\end{frame}

\backmatter

\end{document}
